
\section{Technical Description}

Here you explain the technical work you have carried out. You should present the
overall architecture of your interpreter. 
Try to explain the development choices you made taking into account important properties, 
such as reusability and performance.

\subsection{Architecture}
  
Please provide at least \textbf{one diagram} of the entire 
architecture and describe it. How do the individual elements of the interpreter (i.e., scanner, parser, etc.) work together to execute a given program? Have you applied specific design patterns?

Our interpreter is designed using a multi-stage pipeline architecture. An overview of this pipeline can be seen in figure XX.

At the center of both the figure and the flow of the interpreter is the VVPLController class which orchestrates the entire control flow.
The class steers the entire compilation by invoking the classes responsible for each of the stages of the compilation.
At each stage any discovered errors are reported to the controller which serves as a single centralized error manager
ensuring that the next stage only occurs if the current stage reported no errors. This gurantees that later stages only operate on well-formed inputs
preventing cascading errors.
The use of this design promotes modularity and seperation of concerns, where each component is responsible for a distinct phase of the interpretation
process. The components can focus on their specific responsibility without having to concern themselves with the details of the other components.

Each of the different components surrounding this VVPLController in the figure are labeled with a number
indicating the order of execution.

%Input is VVPL source file.
%Output is the result of program's execution.

\subsection{Implementation}

% If you have implemented specific algorithms or used functionality not covered in the lecture, please explain it here.
% Which kind of design decisions did you take? Answer the following questions in the respective sections.
% \begin{itemize}[noitemsep]
% 	\item How do you recognize tokens in the Scanner?
% 	\item How does the Parser work?
% 	\item How do you perform the semantic analyses?
% 	\item How do you handle the casting of types?
% %	\item How do you (or the grammar) treat the dangling-else problem?
% %	\item Do you allow for syntactic sugar? If so, please report what and where.
% %	\item {What do you check in the scope analysis?}
% %	\item {What do you check in the type analysis?}
% 	\item {How do you handle functions?}
% 	\item What tasks are performed by the interpreter?
% 	\item How and when do you report errors?
% \end{itemize}
\subsubsection{Lexical Analysis}
Scanning is the first step of the pipeline. Here we translate the character sequence from the source program into a token sequence.
We do this by iterating the entire input string trying to recognize different types of tokens.
%! Cuttet fra for plads. Er nok heller ikke så vigtigt?
%Each Token object can store the lexeme, literal, the line the token appears at and the type of the token which is one of the values from enum \texttt{TokenType}.
%To track where in the file it is at, the Scanner uses three pointers: start, current and line.
%The start pointer always indicates the start of the lexeme being scanned while current points to the character
%we are currently inspecting. The line field keeps track of the line we are at in the source program.
In each iteration we scan a single token by utilizing the \texttt{scanToken} method consiting of
one big switch statement to inspect a character and match it to the corresponding case.

The switch statement has a series of simple single character cases which
includes parenthesies, braces, coomas, semicolons, and minuses.
All of these cases are easy - we simply generate the a token with the corresponding type.
Other single character lexemes include comments, newlines and whitespaces.
All of these aren't meaningful to the parser and as such we perform the necessary logic to skip them. 


When encountering longer lexemes like strings, numbers and identifiers we call separate functions to handle these.
For strings we attempt to consume until the closing " is found reporting an error if the string spans too many lines.
When we see a number, we consume the entire series of digits allowing for a decimal point with one or more trailing digits.
%! Kan måske godt bare cuttes til: "leading and trailing decimal points are not allowed".
By looking ahead two characters we ensure a digit is present after the '\texttt{.}' disallowing trailing decimal points.
Leading decimal points are also not allowed, as the character '\texttt{.}' matches to no case in the switch statement.
The literal value is stored as a Double.

Identifiers must start with either a letter or an underscore but can afterwards be followed by digits.
When handling identifiers we start off by consuming the entire literal.
This is in accordance with the maximal munch principal where a token is the longest possible valid string.
The entire character sequence is matched against the map of all reserved keywords.
Only if the sequence does not match any reserved keyword it is categorized as an identifier.
This order of operations is deliberate to ensure that keywords have a higher priority than user-defined identifiers.

%! Kan bare slette dette også. Prøvede at lave en nice overgang, men måske er det ikke nice.
If no errors are detected, we know the input was without lexical errors and is now converted into a 
sequence of tokens that we can use in the next step: Parsing.
\subsubsection{Syntactic Analysis} 
%Because tokens are longest possible valid string and because we first match keywords, then identifiers, true(); is a parse error.
%ikke derfor men ja det er en parse error.

%TODO: Husk at nævn afvigelser fra tests.


% Idéer til afsnit
% 1. Overordnet om parser
%   - performs syntactic analysis of VVPL programs
%   - Uses tokens from scanner to construct an AST
%   - What this AST consists of (Expr, Stmt, polymorphism)
%   - A tree that matches the grammatical structure of the language
%   - A recursive descent parser. 
%   - Noget med at forklare den overordnede rekursive natur og hvordan parser blot mapper tokens til den rekursive grammar
% 2. Error recovery (synchronization)
%   - Tænker dette er noget af det lidt fede ved parseren og god blanding af teori og det tekniske
%   - Kan genbruge mange af idéerne fra monitoring and logging afsnttet, som ikke længere skal bruges.
% 3. Detaljer
%   - Flet nogle af detaljerne ind undervejs. 
%   - Vigtigste er: håndtering af casts, consume rapporterer previous og måske true()



The parser performs syntactic analysis of VVPL programs.

Recursive-descent parser based directly on the VVPL grammar.

Consumes a list of Tokens produced by the Scanner and constructs an abstract syntax tree (AST).
This conists of Expr and Stmt nodes.
We use polymorphism to allow the use of the more detailed subclasses.


Detaljer som måske er værd at nævne:
- Wrapper block stmt. Kan ikke helt huske hvorfor dette var smart, men mener også bogen gør det.
- Håndtering af casts. Vi har valgt at have en seperat CAST node.
- Vi håndterer true() som en parse error selvom det er beskrevet som en scope error i projektbeskrivelsen
- Vi bruger en custom "Param" klasse til at håndterer funktioners parametre.
- Consume rapporterer error på previous().
    Er fint fordi vi altid matcher inden vi kalder consume. Så vi ved altid at mindst én token er
    fundet inden vi kalder consume. Så vi kan ikke komme out of bounds.
    Valget er taget for at manglende semikoloner bliver rapporteret på den rigtige linje.
    Lidt ulogisk i forhold til at vi ikke rapporterer error på den token der faktisk skabte fejlen.
    Men pga. vi kun rapporterer linje, så virker det godt i dette tilfælde.
\subsubsection{Semantic Analysis} 
After having making sure that the program is syntactically correct, we want to confirm that our context-sensitive grammar is also semantically correct, 
i.e. perform context-sensitive checks and interpret the meaning of the program at its' syntactic level.
VVPL is a statically typed language, which usually means that we perform context-sensitive checks at compile-time. 
VVPL is however an interpreted language, meaning that we do not compile the code directly to native machine code, but let the interpreter execute the code.
This means that in our case, "statically-typed" means that we perform the context-sensitive checks before interpreting the program (runtime evaluation).

Context-sensitive checks are performed by constructing symbol tables upon entering a new scope.
VVPL practises lexical (or static) scoping by using curly brackets to indicate a new scope. 
An example of constructing a new symbol table to represent the scope of the current block can be seen in the visitBlockStmt:

\begin{lstlisting}[language=java]
public Void visitBlockStmt(BlockStmt blockStmt) {
        SymbolTable oldTable = currentEnvironment;
        currentEnvironment = new SymbolTable(currentEnvironment);
        for (int i = 0; i < num_statements; i++) {
            analyse(blockStmt.stmts.get(i));
        }
        currentEnvironment = oldTable;
        return null;
}
\end{lstlisting}

Here we construct a new symbol table which represents the scope of the current block. The new symboltable has the old symbol table as an attribute such
that it is possible to access variables from outer scopes, where as variables in inner scopes will be unreachable.

\paragraph{Scope Analysis}\mbox{}\\
Roughly we handle three things in the scopeanalyzer, being variables, functions declarations and function calls. We will go through them step-by-step

\subparagraph{Variables}\mbox{}\\
The following scopeing rules fall in this category (omitting the variables within functions for now):
\begin{itemize}
    \item Variable can only be declared once by a variable declaration
    \item Variables can not be shadowed (i.e. defined in an inner scope to shadow the object in the outer scope)
    \item Variable declarations needs to be initialized with a value
    \item Variable must be visible in current scope before usage
\end{itemize}

scopeanalyzer.java enforces the first two bullet points in the logic of SymbolTable.java. Upon attempting to define, we check if the variable exists in the current or any outer tables to prevent shadowing.

\begin{lstlisting}[language=java]
    public void define(String symbol, Token token) throws SymbolTableException {
        if (contains(symbol)) {
            throw new SymbolTableException();
        }
        symbols.put(symbol, token);
    }
        public boolean contains(String symbol) {
        if (symbols.containsKey(symbol)) {
            return true;
        }
        if (outer != null) {
            return outer.contains(symbol);
        }
        return false;
    }
\end{lstlisting}

We that declarations needs to be initialized by a value in visitVarDecl:
\begin{lstlisting}[language=java]
if (varDecl.expr != null) {
    analyse(varDecl.expr);
}
else {
    VVPLController.error(varDecl.id.line, ErrorTypeStrings.SCOPE_ERROR, "declared variable must be initialized with a value.");
}
\end{lstlisting}

In visitIdentifier we enforce the fourth bullet point by checking if the symbol table contains the identifier:

\begin{lstlisting}[language=java]
public Void visitIdentifierExpr(Identifier identifier) {
    if (!currentEnvironment.contains(identifier.id.lexeme)) {
        VVPLController.error(identifier.id.line, ErrorTypeStrings.SCOPE_ERROR, "variable [insert name here] does not exist in scope or any parent scopes.");
    }
}
\end{lstlisting}

\subparagraph{Function Declarations}\mbox{}\\
The following scopeing rules fall in this category:
\begin{itemize}
    \item Variables inside the function declaration cannot access any outer scopes. These must be \textbf{shadowed}
    \item cannot be defined in inner scopes, only in the global scope
    \item Unique function name (i.e. re-declarations are not possible)
    \item Return statements can only occur in the context of a function
    \item No statements can follow a return statement
    \item Out-of-order declarations and calls are allowed, i.e. a function call can precede its' declaration
\end{itemize}

In visitFunctionStmt, we enforce the first bulletpoint by constructing a new symbol table with no outer symboltable for functions. This way they cannot interact with any values defined outside of this block.
The second bulletpoint is enforced by checking whether the current environment is the original defined environment for the program.
This environment is made unique by adding a separate constructor for a symbolTable marking it as the global environment.
One could think we could check whether our global environment is global by seeing if there are no outer symbol tables linked to the current symboltable,
but due to our solution for encapsulating functions in their own scope, this is not possible.

We have constructed a separate symbol table class to contain functions - this of course has the consequence that a variable and a function can share the same name, but it beautifies the implementation of our scopeanalyzer significantly.
This new table maps a symbol not only to their token, but to their respective FunctionStmt AST node. This way additional details such as the function block and the function type can easily be fetched.
This symbol table also has a much more restricted functionality.
1. It does not need an assign function - we cannot reassign a value to a current existing function
2. It does not need the attribute outer, since the function table exists globally.
The third requirement to function scopes is enforced by this symbol table. 

\begin{lstlisting}[language=java]
    public void define(String symbol, FunctionStmt functionStmt) throws SymbolTableException {
        if (contains(symbol)) {
            throw new SymbolTableException();
        }
        symbols.put(symbol, functionStmt);
    }
    public boolean contains(String symbol) {
        if (symbols.containsKey(symbol)) {
            return true;
        }
        return false;
    }
\end{lstlisting}
In contrast to earlier, we do not need to iterate over any outer symbol tables.

To enforce that return statements can only occur in a function, a global variable is\_function\_env is used.
This variable is set to true upon analyzing the body of a function. This way, return statements will return an error if they occur in a function.
\begin{lstlisting}[language=java]
public Void visitReturnStmt(ReturnStmt returnStmt) {
    if (!is_function_env) {
        VVPLController.error(returnStmt.returnKeyword.line, ErrorTypeStrings.SCOPE_ERROR, "Return statement cannot occur outside of function.");
    }
}
\end{lstlisting}

visitBlockStmt enforces the fifth bullet point by checking that a returnStmt cannot be a non-final statement of its' containing block:
\begin{lstlisting}[language=java]
for (int i = 0; i < num_statements; i++) {
    if (blockStmt.stmts.get(i) instanceof ReturnStmt && i != num_statements - 1 ) {
        VVPLController.error(((ReturnStmt)blockStmt.stmts.get(i)).returnKeyword.line, ErrorTypeStrings.SCOPE_ERROR, "Return statement must be last statement of block.");
        return null;
    }
}
\end{lstlisting}

Lastly, we enforce out-of-order declarations and calls by iterating through the list of given statements and analyse function declarations prior to any other statement:
\begin{lstlisting}[language=java]
public void analyse() {        
    ListIterator<Stmt> stmts = program.listIterator();

    while (stmts.hasNext()) {
        Stmt currentStmt = stmts.next();
        if(currentStmt instanceof FunctionStmt) {
            analyse(currentStmt);
            stmts.remove(); 
        }
    }
    // Analyse rest of program
    for (Stmt stmt : program) {
        analyse(stmt);
    }
    return;
}
\end{lstlisting}

\subparagraph{Function Calls}\mbox{}\\
The following scopeing rules fall in this category:
\begin{itemize}
    \item Only declared and known functions can be called
    \item Number of function arguments matches number of function parameters
    \item Function arguments must be within scope
\end{itemize}
These are all enforced within visitCallExpr. The FuncSymbolTable throws an exception in the case that an unknown function is called.

\begin{lstlisting}[language=java]
    // check if #args == #params
    if (params.size() != args.size()) {
        VVPLController.error(call.id.line, ErrorTypeStrings.SCOPE_ERROR, "function [name here] takes exactly [params.size()] parameters.");
        return null;
    }
    // check if arguments are in scope
    for (int i = 0; i < args.size(); i++) {
        analyse(args.get(i));   
    }
\end{lstlisting}
The given arguments are analysed just like any other expression. This concludes our scope analysis.

\paragraph{Type Analysis}\mbox{}\\
We have constructed a separate file for doing the type analysis. This could be done in the same file as before, but SOLID principles such as the Single Responsibility Principle suggests that responsibilities should be delegated to each their class.
Therefore we build the symbol tables from scratch again in order to build a context-sensitive environment, but we do not need to error handle scope issues since those have been handled already.
Again we can split the typeing requirements in three and go through them step-by-step.

\subparagraph{Variables}\mbox{}\\
The following scopeing rules fall in this category:
\begin{itemize}
    \item Variables are initialized with the correct type
    \item New assignments to variables matches the current type
\end{itemize}

We enforce the initialization of the correct type in varDecl as seen below: 
\begin{lstlisting}[language=java]
if (exprType != castToType) {
            VVPLController.error(varDecl.id.line, ErrorTypeStrings.TYPE_ERROR, String.format("Type %s does not match type of expression (%s)", exprType, castToType));
}
\end{lstlisting}

Likewise, we see if new assignments are compatible in terms of type in visitAssignExpr:
\begin{lstlisting}[language=java]
if (currType == exprType) {
        return currType;   
}
else {
    VVPLController.error(assign.ID.line, ErrorTypeStrings.TYPE_ERROR, "current type of identifier [insert name here] does not match the type of the given expression.");
}
\end{lstlisting}

\subparagraph{Operations}\mbox{}\\
The following scopeing rules fall in this category:
\begin{itemize}
    \item Arithmetic operations only between numbers
    \item Logical Operations only between bools
    \item Conditionals (used by if-statements and while-statements) must be of type Bool.
    \item Only certain type casts are correct
    \begin{itemize}
        \item Number to String
        \item Number to Bool
        \item String to Number
        \item Bool to String
    \end{itemize}
\end{itemize}

Arithmetic Operations are handled in visitBinaryExpr, and it is here required that arithmetic operations must be between numbers.
Likewise, logical operations are handled in visitLogicalExpr.
Requiring conditionals of type bools are done in those functions where AST-nodes with conditions are handled - those are visitIfStmt and visitWhileStmt.
Type rules within casting are evaluated within visitCastExpr.

\subparagraph{Functions}\mbox{}\\
We have made the assumption that if no explicit type has been declared, the function is required to return void.

The following scopeing rules fall in this category:
\begin{itemize}
    \item The types of a function calls given arguments must match the types of the declared parameter(s).
    \item A function declaration must in all cases return the declared type. It cannot miss a return statement or return the wrong type.
\end{itemize}

We enforce the first bullet point in visitCallExpr.
The enforcement of the second bullet point requires more careful handling. This can be divided into two more requirements:
\begin{itemize}
    \item Return statements within a function must all return the declared type
    \item A function must always guarantee to return the correct type nevermind which branch it chooses to take. 
\end{itemize}

We start with the first point. Before visiting any ReturnStmt, the global variable currentFuncReturnType has been set to the desired type in visitFunctionStmt:
\begin{lstlisting}[language=java]
if (functionStmt.typeToken != null) {
    currentFuncReturnType = convertVariableType(functionStmt.typeToken.type);
}
else {
    currentFuncReturnType = Type.UNKNOWN;
}
\end{lstlisting}

In visitReturnStmt we first check if the value of return statement is null. If it is, and the declared type is not null, an error is raised:
\begin{lstlisting}[language=java]
public Void visitReturnStmt(ReturnStmt returnStmt)
    if (returnStmt.returnValue == null) {
        // if current function has defined type then it should not be able to just return with no value
        if (currentFuncReturnType != Type.UNKNOWN) {
            VVPLController.error(returnStmt.returnKeyword.line, ErrorTypeStrings.TYPE_ERROR, "Function with return type must return a value.");
        } else {
            // Void function can return with no value so this is fine.
            return null;
        }
    }
\end{lstlisting}

If the return type of the current function is however not null, we compare it to the declared return value of the function and raise an error if they do not match.
\begin{lstlisting}[language=java]
if (exprType != currentFuncReturnType) {
    VVPLController.error(returnStmt.returnKeyword.line, ErrorTypeStrings.TYPE_ERROR, "Type of returned value does not match declared return type of function");
    return null;
} else {
    return null;
}
\end{lstlisting}

Now for the second point: "A function must always guarantee to return the correct type nevermind which branch it chooses to take". This is done in visitFunctionStmt:
\begin{lstlisting}[language=java]
if (functionStmt.typeToken != null && !alwaysReturns(functionStmt.body)) {
    VVPLController.error(functionStmt.name.line, ErrorTypeStrings.TYPE_ERROR, "Function must return a value in all code paths");
}
\end{lstlisting}







\subsubsection{Tree-Walk Interpretation} 


% • How do you handle the casting of types?
% • How do you handle functions?
% • What tasks are performed by the interpreter?
% • How and when do you report errors?


The interpreter is the section of our program actually evaluating the preprocessed source code and producing the results for the user.
In essence, this is where the program output is produced. 
We are interested in executing the statements and evaluating the expressions. This is done by walking the AST via the \verb|Interpreter| visitor. 
Importantly, this requires an environment that manages variables and the scope dynamically during program execution. This environment can be considered a runtime equivalent of the beforementioned Symbol Table, and is therefore implemented likewise.
This stores our runtime binding of variables and represents the internal state of our interpreter. 
% Måske "when visiting statements" eller "executing statements" i stedet for bare "the statements"? - Lucas
The statements changes the internal state and the expressions evalautes to an \verb|Object|.
% As the actual value is only later evaluated? Måske er noget i den dur bedre? - Lucas
\verb|Object| is an important runtime artifact, as we need to evaluate the actual value later when interpreting statements. 

% The fact that in our implementation of VVPL, the semantic analysis has already handled scope and type checking leaves us with much fewer concerns regarding error handling in this section. 
% LUCAS:
% Ville ikke skrive "consulting with Sandra". Synes jeg er lidt mærkeligt. 
% Måske i stedet brug pladsen på at forklare lidt tydeligere problemet: Når vi skal caste fra streng til tal bliver vi nødt til at kende hvad strengens værdi ender med at blive evalueret til.
% Kan også godt være man forstår det fint bare af at du skriver: Some errors can only be detected at runtime; for example errors relating to the actual literal value of variables after being evaluated.
Despite the semantic analysis, there are some unique error concerns for our interpreter. We need to abort the program immediatly upon a runtime error. 
Some errors can only be detected at runtime; for example errors relating to the actual literal value of variables after being evaluated. We have decided to implement the handling of one type of error, namely the conversion of \verb|String| to \verb|Number| after consultation with Sandra.
A small note is that dividing by zero is handled per default in Java, but refering to the project description, we can in general assume no runtime exceptions. We do a test for this, just to document the behaviour. 

Functions are defined via a \verb|VVPLCallable| interface, and implemented via the \verb|VVPLFunction| class. 
A benefit is future extension; having defined an interface allows us to more easily extend our language with for example built-in functions. 
% Hvad betyder det helt præcist at vi håndterer hver funktion som sin egen instans? Det der står lige efter? 
% At vi har et nyt environment? Måske mere nysgerrighed end en kommentar - Lucas
We then handle each function call as its own instance; importantly, we pass on the provided arguments and the global variables in a new environment called \verb|functionScope|. This is what allows us to modify a copy of the global variables, and maintain the \textit{call by value} porperty of the functions.
Within this scope, we can now execute the block body statements of the function. Notice that we already check for the arity of the function parameters and provided argument in the \verb|ScopeAnalyzer| visitor, and for types of the function parameters and prodivded arguments in \verb|typeAnalyser|, so these checks are not performed.
The function call can consist of an arbitrary amount of visited statements.
If there is a value to return after the executed statements of the function, we must get back to the top of the call stack with the return value. 
This is implemented with the \verb|ReturnExcep| class, whose instance stores the return value as an attribute. This instance is then thrown in the \verb|visitReturnStmt| method, and caught and returned by the top-level \verb|visitCallExpr| function call, and the program can now carry on executing statements. 

% Lucas:
% Ah okay du forklarer tydeligt cast problemet her. Så ville jeg bare slet ikke nævne det tidliger.
% Altså ville slette denne tidligere sætning helt: We have decided to implement the handling of one type of error, namely the conversion of \verb|String| to \verb|Number| after consultation with Sandra.
Casting of types is handled by the \verb|visitCastExpr| method. The casting of a \verb|String| to a \verb|Number| requires us to be aware of the actual value of the string, which is only available at runtime. This is implemented by trying to convert the string to a \verb|Double| with the built-in java method \verb|parseDouble|. If this do not succed, we know the content of the String is not valid, and we can catch the runtime error.
Another notable casting feature, is the fact that all numbers should evaluate to \verb|False| except for the number 13. This is implemented inside of the helper function \verb|isTruthy|, where a hard-coded check for the number 13 is conducted.

% \textbf{Testing the interpreter}


\subsubsection{Monitoring and Logging} 
%Do you trace the state of the interpreter and its components (i.e., the scanner, parser, etc)?
% % Explain how we do error handling.

An essential part of making a programming language usable includes having a sensible error handling.
Clear error messages help users understand what went wrong in their program and how to fix it.
To facilitate this, we perform error handling in all parts of the interpreter.

At the center of the error handling is the \texttt{VVPLController} class which functions as a centralized hub for
reporting errors. This design decouples the error \emph{detection} from the error \emph{reporting}.
Each component such as the scanner or parser is responsible for identifying errors and notifying the
\texttt{VVPLController} class which in turn is responsible for handling how these errors are stored
and reported to the user.

This interplay between the classes is facilitated by the \texttt{error} method in the controller.
\begin{lstlisting}[language=java]
public static void error(int line, String errorType, String message) {
    hadError = true;
    errorMessages.add(new ErrorMessage(line, errorType, message));
}
\end{lstlisting}

This method takes in the line where the error occured, the type of error and the error message we are reporting to the user. 
The error type is one of the different errors from the \texttt{ErrorTypeStrings} class, which allows us to categorize the errors.
The method firstly sets the \texttt{hadError} flag, which signals that the current stage of compilation has failed.
After each stage of the interpreter pipeline, the controller checks this flag. If an error has been registered the pipeline stops.
Halting the interpreter as soon as a single stage fails is part of a fail fast strategy that prevents later stages from operating 
on invalid data which would cause a cascade of confusing and possibly misleading errors.
%TODO: Skal denne sætning om at scanner, parser etc. reporter alle errors. Kan måske i stedet bruges senere til naturlig overgang.
While the pipeline stops between stages, each individual stage is designed to detect as many errors as possible before terminating,
providing the user with a comprehensive report of their errors in that stage.
This gives the user a more complete overview of issues which decreases time spent debugging by
reducing the number of compiles needed to find and fix all errors in the program.

After setting the \texttt{hadError} flag the error is added to the errorMessages list which stores all the errors.
This is a list of \texttt{ErrorMessage} objects:
\begin{lstlisting}[language=java]
private static List<ErrorMessage> errorMessages = new ArrayList<>();
\end{lstlisting}

The \texttt{ErrorMessage} class stores errors as structured data rather than just plain strings.
It stores the line number, error type and message in seperate fields which allows for easy post-processing.
Specifically this allows for sorting the 
errors by line number such that the user receives a logical and actionable list of issues in the order they appear in the code. 
This is achieved by having \texttt{ErrorMessage} implement the \texttt{Comparable} interface allowing the controller to sort the list with
the \texttt{sort} method from the Collections class.
The class also overrides the \texttt{toString} method such that the error messages can easily be reported to the
user in the format: \texttt{<error-type>, line <lineNo> <message>}.

% Kan måske passe ind et sted?:
% Our system monitors the state of the source code's correctness at each stage and logs deviations in a structured user-friendly manner.



\subsection{Testing \& Evaluation}

How did you test your interpreter? Which coverage criterion did you apply?\\
Describe the test setup and the test inputs with their expected outputs.
If you have documented this in the source code a link and submission of the generated JavaDoc is also fine.

If the provided tests have been changed or do not work, please explain here what tests fail, why and what it would require to fix them.

Overall, our testing is characterized by both testing on a smaller and bigger scale. We test the individual ... throughout all stages, and then test whole end-to-end pipeline with the interpreter in \verb|VVPLController|

% Idé til struktur?: /Lasse
% - Unit test (vi tester de enkelte dele hver for sig, e.g expressions i en parser)
% - Integration test (Vi tester, at scanner + parser virker sammen og printede det i AST (del 1))
% - end-to-end test (Vi tester den fulde pipeline med tilhørende interpretation, hvor vi sammenligner source code med evaluated code)

% Overvej dog om Sandra egentlig går op i det. har bare skrevet en lille kommentar om det, måske slet. 

% Herefter en kort kommentar omkring vores code coverage. 

% Synes igen det er ret mærkeligt at nævne Sandra ved navn.
% Gode tanker.
% Når vi når dertil ville jeg ændrer det til et mere overordnet blik
% Og derfor er opremsningen af de ting vi testede lige i interpreter nok ikke så vigtigt.
% Men tænker vi siger vi også havde et fokus på edge cases vi fint kan komme med f.eks. RuntimeException eksemplet.

\begin{comment}
How did you test your interpreter? Which coverage criterion did you apply?\\
Describe the test setup and the test inputs with their expected outputs.
If you have documented this in the source code a link and submission of the generated JavaDoc is also fine.

If the provided tests have been changed or do not work, please explain here what tests fail, why and what it would require to fix them.
\end{comment}


Our testing strategy focused on complementing the provided test cases with tests that exercise different branches and edge cases.
The scanner, parser and semantic analysis is mainly tested using negative tests. That is tests with invalid input that ensure
errors are detected and handled gracefully.
In contrast, the testing of the interpretation process is primarily positive tests that confirm
programs are evaluated as expected. 
Additionally we made sure to test the \verb|RuntimeException| behavior when casting an invalid string to a Number to check 
if it interrupts the program flow as intended.
These interpreter tests also serve as end-to-end tests of the entire pipeline, 
since correct interpretation depends on the correctness of all preceding stages.

To evaluate our tests we generated a coverage report, which helped identify previously untested cases.
With the addition of our test cases an instruction coverage of 96\% and a branch coverage of 88\% was reached.
The remaining uncovered branches primarily stem from unreachable code and simple checks for cascading errors. 
As such we consider the reached coverage to be sufficient.

The provided \texttt{parse-errors.in} test fails, as the error on line 12 is not reported.
The relevant sections of the test can be seen in the snippet below.
This occurs because the missing semicolon on line 10 causes the parser to enter panic mode.
During synchronization we deliberately do not synchronize on identifiers, since identifiers
do not gurantee the start of a new statement. As a result the entire assignment on line 12
is discarded until the semicolon at the end of that line is consumed, which ends the synchronization.
The problem stems from the exprStmt grammar rule through which expressions can become statements.
So syntactically an assignment is not a statement itself but becomes a statement when followed by a semicolon.

\begin{lstlisting}[language=Java, firstnumber=10]
variable a2 of_type Bool is true 	# missing semicolon

a is add(b,c);

if ( ) 							#condition missing
	a is add(b,c);
\end{lstlisting}

Adding identifiers as synchronization points would of course not work, as identifiers appear in the middle
of all sorts of expressions, leading to cascaded errors if we were to synchronize on these.
We could specifically detect assignment expressions by checking whether an identifier is followed by a token of type assignment.
However, since assignments can also appear inside expressions this would only be a recovery heuristic and would break the current invariant 
that synchronization occurs only at grammar-guranteed statement boundaries. 
Futhermore, it would still not make the test pass, as with this strategy we would also report the similar error on line 15,
which is not reported in the expected output.
We decided to stick with a simpler approach where the synchronise method returns when it is a grammar gurantee that a new statement has started.
The test itself seems to be inconsistent with whether we should synchronize on assignments or not as catching the error on line 12 but not on line 15 indicates.
Although our recovery strategy may be conservative it ensures we don't resume inside broken expressions.


The \texttt{parse-errors.in} test also shows a minor discrepancy in that the missing statement on line 20 
is reported on line 21 instead.
This occurs because the parser reports the error when encountering the unexpected \texttt{else} token where a statement was expected.
We consider this acceptable, as the \texttt{then} statement could equivalenty have appeared on line 21, making the reported error both valid and understandable.

The sample-ast-expected.out file fails on the final 3 lines, since we have a dedicated node for casts in the AST.
The test would pass by adjusting the expected output to reflect this structure, specifically by simply
swapping the order of the \texttt{LiteralExpr} and \texttt{Cast\_To} nodes with appropriate indentation.

\begin{comment}

Overall, our testing is characterized by both testing on a smaller and bigger scale. We test the individual ... 
throughout all stages, and then test whole end-to-end pipeline with the interpreter in \verb|VVPLController|

For the interpreter, our additional tests is oriented towards testing cases not already taken care of by the other stages, 
and not already part of Sandra's test suite. This encompass the \textit{call by value}, casting from a Bool to a String, 
casting numbers to Bools (including 13), casting Strings to Numbers, division by zero and finally, the remaining untested 
expression in the form of the \verb|GREATER| binary expression.
Importantly, we made sure to test the \verb|RuntimeException| behavior when casting an invalid string to a Number to check 
if it interrupts the program flow as intended. The division by zero was simply to document the behaviour, as we relies on 
Javas default error handling of this. 
\end{comment}


%1. overall teststrategy
%2. coverage
%3. failed tests (2x in parser, 1x in AST printing)



